\documentclass[12pt]{article}
\usepackage[utf8]{inputenc}
\usepackage[T1]{fontenc}
\usepackage{hyperref}
\title{Suganga Language Resources for K-12 Curriculum Development}
\author{Compiled by Claude Sonnet 5.5 (claude-sonnet-5-5) for Andr\'as Kornai}
\date{}
\begin{document}
\maketitle

\section{Introduction}
Suganga (ISO 639-3: sug; Glottocode suga1249) is spoken in Papua New Guinea. The speakers live in the border area between East Sepik (Ambunti-Dreikikir) and Sandaun (Telefomin), near the Sepik headwaters. Glottolog places it in the Ok branch (Mountain Ok, Mianic). Joshua Project calls it a dialect known as West Mian or Wagarabai and gives a population of about 600--700. A PARADISEC entry describes a recording from ``Sugan'' as a dialect of Yau, so the classification is uncertain. The language is listed as endangered.

\section{Primary Educational Resources}
None found.

\section{Linguistic and Reference Materials}
\begin{itemize}
\item Glottolog language page (map view): \url{https://glottolog.org/resource/languoid/id/suga1249.bigmap.html}
\item Endangered Languages Project bibliography: \url{https://endangeredlanguages.com/lang/10676/bibliography?hl=fr}
\end{itemize}

\section{Religious and Community Resources}
Joshua Project reports no Scripture translation: \url{https://joshuaproject.net/people_groups/15107/PP} and \url{https://www.joshuaproject.net/index.php/languages/sug}

\section{Digital and Audio Resources}
A PARADISEC catalogue entry describes a recording made in ``Sugan''. It labels the language a dialect of Yau, so it may not be Suganga: \url{https://catalog.paradisec.org.au/repository/VG3/140421}

\section{Curriculum Development Context}
Little exists. A curriculum would need new documentation and orthography work with the community, starting with the PARADISEC material if it is confirmed to be Suganga.
\end{document}